\documentclass[FGBook.tex]{subfiles}

\begin{document}

\drama{Revizor}{NIKOLAJ VASILJEVIČ GOGOL}

\noindent\begin{wrapfigure}{r}{0.35\textwidth}\footnotesize
\subwection{Ukázka z díla:}
\setlength{\parindent}{3pt}
HEJTMAN: Pánové, přijel revizor! Inkognito! A ještě s tajným příkazem!\\
BOBČINSKIJ a DOBČINSKIJ (současně): Ano, ano! Viděli jsme ho v hospodě!\\
HEJTMAN: Co teď? Musíme všechno uvést do pořádku... rychle!
\end{wrapfigure}

\wection{LITERÁRNÍ TEORIE}

\subwection{Literární druh a žánr:}
\noindent drama, satirická komedie (komedie omylů s groteskními prvky)

\subwection{Literární směr:}
\noindent Kritický realismus (zakladatel v ruské literatuře). Navazuje na romantismus, ale přechází k realistickému obrazu společnosti s prvky grotesky a satiry.

\subwection{Jazyk a slovní zásoba:}
\noindent Hovorový, živý jazyk s vulgarismy, archaismy a typickými ruskými příjmeními (mluvící jména – Skvoznik-Dmuchanovskij, Chlestakov). Dialogy odhalují charaktery – hejtman mluví důstojně, Chlestakov nafoukaně a nesouvisle.

\subwection{Figury:}
\noindent Ironie, opakování, kontrast (zdánlivá důstojnost vs. korupce), groteska.

\subwection{Tropy:}
\noindent Satira, alegorie (městečko = celé Rusko), symbolismus (revizor = svědomí / Soudný den), hyperbola.

\subwection{Stylistická charakteristika textu:}
\noindent Pět jednání s chronologickým postupem. Rychlé dialogy, scénické poznámky, groteskní situace. Smích z absurdity a lidské malosti („Nevrč, brachu, na zrcadlo, když máš křivou hubu“).

\subwection{Postavy:}
\noindent 
\textbfhl{IVAN ALEXANDROVIČ CHLESTAKOV --} lehkomyslný, přitroublý petrohradský úředníček, povaleč a lhář. Využívá omylu a užívá si roli „revizora“.\\
\textbfhl{ANTON ANTONOVIČ (HEJTMAN) --} korupčník, lstivý, ale zbabělý. Bojí se kontroly, podplácí a lichotí.\\
\textbfhl{ANNA ANDREJEVNA a MARIE ANTONOVNA --} hejtmanova žena a dcera – marnivé, koketní, snící o Petrohradu.\\
\textbfhl{BOBČINSKIJ a DOBČINSKIJ --} zvědaví, hloupí zemani, kteří celý omyl odstartují.\\
\textbfhl{OSIP --} Chlestakův sluha, chytřejší než pán, pragmatický.

\subwection{Děj:}
\noindent V malém ruském městečku se rozšíří zpráva, že přijede tajný revizor. Zkorumpovaní úředníci (hejtman, soudce, školní inspektor, poštmistr aj.) se v panice snaží zakrýt nepořádek a úplatky. Dva zvědaví zemani považují za revizora lehkomyslného petrohradského úředníčka Chlestakova, který v hospodě nemá na zaplacení. Hejtman ho pozve k sobě domů, všichni mu lichotí a půjčují peníze. Chlestakov se vžije do role, naparuje se, dvoří se hejtmanově ženě i dceři a slibuje sňatek. Nakonec odjede. Poštmistr otevře jeho dopis a všichni zjistí, že to byl podvod. Závěr: přijíždí pravý revizor.

\subwection{Kompozice, prostor a čas:}
\noindent Pět jednání, chronologický postup (děj proběhne během cca 24 hodin). Prostor: malé provinční ruské městečko (neurčené). Čas: 30. léta 19. století (vláda cara Mikuláše I.).

\subwection{Význam sdělení:}
\noindent Ostrá satira na korupci, byrokracii, maloměšťáctví a lidskou hloupost v carském Rusku. Motto „Nevrč na zrcadlo, když máš křivou hubu“ – kritika celé společnosti. Gogol ukazuje, jak strach z moci vede k podlézavosti a pokrytectví. Alegoricky: městečko = celé Rusko, Chlestakov = falešná autorita.

\vspace*{2em}\\
\wection{LITERÁRNÍ HISTORIE}

\subwection{Politická situace:}
\noindent 30. léta 19. století – vláda cara Mikuláše I. (autokratický režim, cenzura, policejní dohled). Gogol kritizuje provinční korupci a byrokracii.

\subwection{Kontext literárního vývoje:}
\noindent Přechod od romantismu ke kritickému realismu. Námět údajně pochází od Puškina. Gogol ovlivnil Dostojevského, Čechova a celou ruskou literaturu.

\subwection{Nikolaj Vasiljevič Gogol {\ssmall -- život autora:}}
\noindent (1809–1852). Ukrajinského původu, ruský prozaik a dramatik. Studoval v Nižynu, žil v Petrohradě a Římě. Trápila ho deprese a náboženské krize. Zemřel v 42 letech po dlouhém půstu.

\subwection{Další autorova tvorba:}
\noindent Mezi nejznámější díla patří \textsc{Mrtvé duše}, \textsc{Večery na samotě u Dikanky}, \textsc{Taras Bulba}, povídky jako \textsc{Plášť} nebo \textsc Nos}.

\subwection{Aktuálnost tématu a zpracování tématu:}
\noindent Nadčasová kritika korupce, byrokracie, podlézavosti a strachu z moci. Platí i dnes – „revizor“ může přijít kdykoli.

\vspace*{2em}\\
\wection{OSOBNÍ NÁZOR}
\noindent Revizor je geniální satirická komedie – směješ se, ale mrazí tě z toho, jak přesně Gogol vystihl lidskou malost. Chlestakov jako nafoukaný povaleč, který využije situaci, a hejtman, který se z paniky promění v podlézavce, jsou dokonalé typy. Nejsilnější je závěr s pravým revizorem – zrcadlo, které nikdo nechce vidět. Gogol dokázal kritizovat celou společnost vtipně a bez kázání. Rozhodně patří k vrcholům ruské dramatiky – čte se skvěle i dnes.

\vfill

\noindent\begin{minipage}{\textwidth}
    {\textcolor{\wpagecolor}{\rule{\linewidth}{0.4pt}}
    \footnotesize
    \textbfhl{Satirická komedie --} zesměšňuje společenské nedostatky (korupce, byrokracie); 
    \textbfhl{Groteska --} přehánění a zkreslení reality pro komický efekt; 
    \textbfhl{Kritický realismus --} věrný, kritický obraz společnosti bez idealizace.
    }
\end{minipage}

\newpage

\end{document}